\documentclass[conference]{IEEEtran}
\usepackage{cite}
\usepackage{amsmath,amssymb,amsfonts}
\usepackage{algorithmic}
\usepackage{graphicx}
\usepackage{textcomp}
\usepackage{booktabs}
\usepackage{enumitem}

\begin{document}

\title{Conference Paper Title on Academic Publishing Engineering}

\author{\IEEEauthorblockN{First Author}
\IEEEauthorblockA{\textit{Department of Computer Science} \\
\textit{University Name}\\
City, Country \\
email@university.edu}
\and
\IEEEauthorblockN{Second Author}
\IEEEauthorblockA{\textit{Department of Electrical Engineering} \\
\textit{University Name}\\
City, Country \\
email@university.edu}
}

\maketitle

\begin{abstract}
This paper presents the foundational principles and experimental results for modern engineering systems. A comprehensive analysis is provided with mathematical proofs and performance benchmarks.
\end{abstract}

\begin{IEEEkeywords}
Academic publishing, LaTeX automation, compiler optimization.
\end{IEEEkeywords}

\section{Introduction}
Modern scientific communication requires rigorous formatting and reproducible mathematical modeling \cite{shannon1948}.

\begin{equation}
  C = B \log_2\left(1 + \frac{S}{N}\right)
  \label{eq:capacity}
\end{equation}

Equation \eqref{eq:capacity} expresses the Shannon-Hartley channel capacity theorem.

\section{Methodology}
Experiments were conducted using standard peer-reviewed callsets.

\begin{table}[htbp]
  \centering
  \caption{Benchmark Convergence Rates}
  \label{tab:results}
  \begin{tabular}{llr}
    \toprule
    Model & Precision & Latency (ms) \\
    \midrule
    Baseline & 92.4\% & 14.2 \\
    Proposed & 98.7\% & 3.1 \\
    \bottomrule
  \end{tabular}
\end{table}

\section{Conclusion}
The proposed framework achieves superior convergence across all tested configurations.

\bibliographystyle{IEEEtran}
\bibliography{references}

\end{document}
